\begin{center}
[2.]\par\medskip
\textsc{Gauss} an \textsc{Bessel}.\quad Göttingen, 26.~December 1821.
\end{center}

\par\noindent\hspace*{\parindent}.\ .\ .\ .\ .\ .\quad Der vorige Sommer ist grösstentheils mit Vorbereitungen zum Trianguliren und dem Trianguliren selbst zugebracht.\quad.\ .\ .\ .\quad Die Winkel in Sternwarte, Meridianzeichen, Hohehagen, Hils, Brocken sind gemessen; nur am letzten Punkte ist der Winkel zwischen \textls[120]{Inselsberg} und andern Punkten missglückt, weil das unerhört schlechte Wetter während der Zeit, wo dort der Heliotrop war, fast alles Beobachten untersagte.
Nur einmal auf eine halbe Stunde konnte die Richtung kümmerlich gemessen werden, was aber mit andern frühern ebenso kümmerlichen Messungen, wo einmal auf das Haus, das andere Mal auf \textsc{Enckes} Sextanten-Viceheliotrop 3 Minuten vor Sonnenuntergang pointirt wurde, nicht gut harmonirt.

\par\noindent\hspace*{\parindent}Der Umstand, dass ich nur Einen Heliotrop zu meiner Disposition hatte, hielt ungemein auf.
Bei meinem Aufenthalt auf dem Hils musste der Heliotrop successive von Lichtenberg zum Meridianzeichen, dem Brelingerberg und Deister übergehen; ebenso wie ich auf dem Brocken war, reiste jener von Lichtenberg zum Hils, Hohehagen, Inselsberg.
Die wenigsten der Winkel haben also unmittelbar gemessen werden können; ich hatte auf diesen beiden
\par\clearpage
\noindent350\hfill{\small\textbf{BRIEFWECHSEL MIT BESSEL.}}\par\medskip
\noindent Stationen mehrere andere an sich sichtbare Punkte ausgewählt und maass allezeit, was sich eben messen liess.
Ich sehe aus dem 4.~Bande der Base du système métrique, dass die Franzosen in Spanien es ebenso gemacht haben; sie haben aber die Messungen nicht richtig combinirt.
Wenn ich im nächsten Jahre meine Triangulirung fortsetze, wird es besser gehen, indem ich zwei Heliotrope mehr haben werde.
Meinen zum Viceheliotrop eingerichteten Sextanten habe ich immer bei mir geführt und zu telegraphischen Zeichen gebraucht.
Es geht damit ganz vortrefflich, nur dass, weil jener Viceheliotrop auf Winkelabstände unter etwa \(138^{0}\) begrenzt ist, ich nicht immer zu jeder Tageszeit ihn brauchen konnte; künftiges Jahr würde ich zu diesem Zweck in der Regel einen der Heliotrope bei mir behalten.

\par\noindent\hspace*{\parindent}Die grösste Entfernung, wo das Heliotroplicht mit blossen Augen gesehen, ist die vom Brocken zum Hohehagen, indem am letztern Orte meine auf dem Brocken gegebenen telegraphischen Zeichen so gesehen wurden; die Entfernung ist \(9\frac{1}{3}\) geographische Meilen; es gehören dazu aber wohl günstige Umstände.
Bei Entfernungen bis 6 Meilen habe ich (trotz des Verlustes an Licht beim Gebrauch der Lorgnette), wenn die Umstände nur leidlich günstig waren, das Heliotroplicht mit blossem Auge bequem gesehen, ebenso in den Vormittagsstunden das Heliotroplicht vom Hils zum Brocken (\(7\frac{1}{2}\) Meilen).
Die Spiegelfläche ist \(2\frac{1}{2}\) Quadratzoll, bei den beiden neuen Heliotropen habe ich \(6\frac{1}{4}\) Quadratzoll genommen; mässigen kann man das Licht leicht, bei grossen Entfernungen kann es aber doch Fälle geben, wo das stärkere Licht angenehm ist; so drang einen Tag das Licht vom Brelingerberg zum Hils nur kurze Zeit durch, \textls[120]{obgleich} (oder vielmehr richtiger \textls[120]{weil}) ein völlig wolkenloser Himmel war.
An solchen Tagen ist in der Regel die Durchsichtigkeit der Atmosphäre am geringsten; ich erinnere mich eines solchen Tages, wo das Heliotroplicht vom Deister zwar recht gut zu sehen war, aber nicht benutzt werden konnte, weil keiner der übrigen Gegenstände, selbst nicht einmal das nur wenig über eine Meile entfernte Einbeck gesehen werden konnte.

\par\noindent\hspace*{\parindent}Einen sehr wesentlichen Vorzug hat das Heliotroplicht vor jedem andern Signale, worüber ich oft artige Erfahrungen gemacht habe; nemlich das Heliotroplicht sieht man desto besser, je stärker man vergrössert, irdische Signale hingegen (bei grossen Entfernungen) desto schlechter, denn bei letztern ist es vorzüglich die \textls[120]{Blässe}, die das Sehen hindert; von Hannover habe ich
\par\clearpage
\noindent{\small\textbf{ZUR HANNOVERSCHEN TRIANGULATION.}}\hfill351\par\medskip
\noindent z.~B.\ auf dem Brocken, obgleich ich den Platz nach Höhe und Azimuth genau wusste, niemals eine Spur sehen können.
An dem Tage, wo ich das kümmerliche Heliotroplicht vom Inselsberg erhielt (\textls[120]{bloss} Schuld des Wetters, denn bei einigermassen günstiger Luft musste das Heliotroplicht von \(2\frac{1}{2}\) Quadratzoll Fläche noch überreichlich hell sein), sah ich dies zwar im Fernrohr des Theodolithen noch ziemlich gut, konnte damit aber keine Spur vom Umriss des Berges erkennen; dagegen sah mein Gehülfe diesen Umriss ganz leidlich in einem an sich sehr elenden Fernrohr von schwacher Vergrösserung, konnte damit aber das Heliotroplicht nur selten sehen.
Die Sache erklärt sich leicht; auch bei der stärksten Vergrösserung bleibt das Heliotroplicht ein Punkt und dessen Licht ist immer dasselbe, aber der Grund ist desto düsterer, je stärker man vergrössert etc.

\par\noindent\hspace*{\parindent}Grosse Signale habe ich nur zwei gebaut, auf dem Hohehagen und Hils, und vermuthlich wäre auch dies unterblieben, wenn mein Heliotrop sechs Wochen früher vollendet gewesen wäre.
Das Hils-Signal projicirt sich vom Brocken aus gegen nahen dunkeln Hintergrund; während des ganzen Monats, den ich auf dem Brocken zubrachte, habe ich jenes nur zweimal überhaupt \textls[120]{sehen} und nur auf etwa zwei Minuten so sehen können, dass ein Winkel sich hätte messen lassen; das Hohehagen-Signal projicirt sich gegen die entferntern Casselschen Berge, war öfters zu sehen und auch ein paar Mal zu beobachten, obwohl ich, hätte ich nicht den Heliotrop dahin geschickt, auch meine Winkelmessung nicht voll bekommen hätte.
Ausser der schweren Sichtbarkeit, Kosten und Zeitaufwand haben die Signalthürme noch eine sehr unglückliche Seite, den Reiz, welchen sie dem rohen Muthwillen zur Zerstörung darbieten.
Leider ist mein Hohehagen-Signal seit kurzem fast ganz verwüstet, und ich werde glücklich sein, wenn ich nur den Punkt mit hinreichender Schärfe wiederfinden kann.

